\section{Introduction}

A norm may force the average distance to two opposite points to grow in
every sufficiently large direction, yet have no equivalent norm with the
corresponding one-sided growth. We construct two explicit real $L^1$
spaces that exhibit this distinction. In both spaces, multiplying along a
tree has two effects: fresh variables force midpoint growth, while path
products of $L^1$ norm one prevent asymptotically uniformly convex
renormings.

For a norm $N$ on an infinite-dimensional real Banach space $X$, let
$\cof(X)$ denote the closed linear subspaces of finite codimension. For
$t>0$, define the averaged asymptotic midpoint modulus and the usual
one-sided asymptotic modulus \cite[Section~2.1]{JLPS2002} by
\begin{align*}
 \whdelta_N(t)
 &=\inf_{N(x)=1}\sup_{F\in\cof(X)}
   \inf_{\substack{y\in F\\N(y)\ge1}}
   \left(\frac{N(x+ty)+N(x-ty)}2-1\right),\\
 \overline\delta_N(t)
 &=\inf_{N(x)=1}\sup_{F\in\cof(X)}
   \inf_{\substack{y\in F\\N(y)=1}}\bigl(N(x+ty)-1\bigr).
\end{align*}
The norm is \emph{asymptotically uniformly convex} (AUC) if
$\overline\delta_N(t)>0$ for every $t>0$. The subspace $F$ may depend
on the center $x$ and on $t$; it must work for all the indicated directions.
The standard asymptotic midpoint modulus uses unit directions $N(y)=1$
and replaces the displayed average by the maximum of the two norms.
A norm is \emph{asymptotically midpoint
uniformly convex} (AMUC) when this maximum modulus is positive for every
$t>0$ \cite[Definition~2.2]{DKRRZ2016}. All numerical midpoint estimates
below concern the displayed average. They imply the same lower bounds
for the maximum, and hence AMUC, because the maximum dominates the average.

Let $\mathcal T=\N^{<\omega}$ be the tree of finite sequences of positive
integers, with root $\varnothing$. For a nonroot vertex $v$, write $v^-$
for its parent. Fix a positive nonconstant real random variable $W$ such
that
\[
 \E W=1,\qquad m_2:=\E W^2<\infty.
\]
On a countable product probability space take independent copies $W_v$,
indexed by the nonroot vertices, and define
\[
 P_{\varnothing}=1,\qquad P_v=P_{v^-}W_v,
 \qquad X_W=
 \overline{\operatorname{span}_{\R}\{P_v:v\in\mathcal T\}}^{\,L^1}.
\]
We use the inherited real $L^1$ norm, denoted by $\norm{\cdot}$.
Write $X_{\mathrm{exp}}$ for the law with density $e^{-w}$ on $(0,\infty)$,
and $X_{\mathrm{sq}}$ for the law of $G^2$, where $G$ is a standard real
Gaussian. For $K>0$ set
\[
 c(K)=\frac1{16}\E(|G|-2K)_+,\qquad r_+=\max\{r,0\}.
\]

\begin{theorem}\label{thm:main}
For every multiplier $W$ as above, the Banach space $X_W$ is infinite
dimensional and has no equivalent AUC norm.
For the exponential and Gaussian-square laws, the inherited norms satisfy
\begin{align}
 \whdelta_{X_{\mathrm{exp}}}(t)
 &\ge \frac{7t}{30}c(10/t)>0 &&(t>0),\label{eq:main-exp}\\
 \whdelta_{X_{\mathrm{sq}}}(t)
 &\ge\frac{\PP\bigl(|G|\ge4\sqrt2(2+\sqrt2)/t\bigr)}{8}>0
 &&(t>0),\label{eq:main-sq24}\\
 \whdelta_{X_{\mathrm{sq}}}(t)
 &\ge\frac{\PP\bigl(|G|\ge8\sqrt6/t\bigr)}{4}>0
 &&(t>0).\label{eq:main-sq12}
\end{align}
\end{theorem}

The three estimates come from different scalar arguments, whose
parameter-dependent forms are proved below. None of the displayed constants is
asserted to be optimal. The general multiplier hypothesis is used for the
renorming obstruction; the positive midpoint estimates are proved for the
two specified laws.

\subsection*{Context and method}
Dilworth, Kutzarova, Randrianarivony, Revalski, and Zhivkov introduced
AMUC and distinguished it from AUC for a given norm: they constructed an
equivalent norm on $\ell_2$ that is AMUC but is not AUC
\cite[Theorem~2.4]{DKRRZ2016}. The renorming question is different. They
proved that an AMUC space with an unconditional Schauder basis admits an
equivalent AUC norm, and asked whether this remains true for every AMUC
space \cite[Section~5 and Proposition~5.1]{DKRRZ2016}.

Baudier answered that question negatively. He proved that a closed
subspace of $L^1$ whose unit ball is relatively compact in measure is
AMUC, and applied this result to the Kadets--Werner space, which admits
no equivalent AUC norm \cite[Theorem~1 and Corollary~1]{Baudier2026}.
His example also satisfies a linear lower bound for the maximum midpoint
modulus \cite[final paragraph]{Baudier2026}. Theorem~\ref{thm:main}
instead gives explicit spaces obtained from one repeated multiplier law,
with direct proofs of quantitative bounds for the averaged modulus and
of the renorming obstruction.

The construction has a useful relation to the multiplication step in
the Bourgain--Rosenthal construction \cite{BourgainRosenthal1980} as
developed by Kadets and Werner \cite[Lemmas~2.2 and~2.4]{KW2004}.
Their enlargement multiplies selected old vectors by fresh nonnegative
variables, chosen anew at successive enlargement stages. Here one fixed
positive law with finite second moment is repeated at every vertex.
Our arguments concern this closed product span directly; they
make no assertion of reflexivity, measure precompactness of its unit ball,
or the Daugavet property.

The main difficulty in the midpoint estimate is uniformity over all
directions in a subspace of finite codimension. Conditioning on the
variables in a finite predecessor-closed subtree $T$ produces a
finite-rank projection on $X_W$. In a finite linear combination of products
in its kernel, the terms beyond $T$ can be grouped by their first vertices
outside $T$. Each group factors into the multiplier at its first exit and
a coefficient depending only on old coordinates and coordinates strictly
below that exit. Fixing these coefficients leaves the exit multipliers
independent.
An independent-copy and random-sign argument controls the vector's
$L^1$ norm by the expected Euclidean size of its coefficient list.
The scalar probability estimates then force a definite midpoint gain
for every such list, even in the presence of an arbitrary shift.

The three scalar mechanisms are elementary but distinct. For exponential
multipliers, a difference of independent copies is a Gaussian scale
mixture, which gives positive expected excess past every fixed threshold.
For Gaussian squares, an orthogonal change of variables turns the
difference into a conditionally Gaussian variable. One proof transfers
its tail probability back to the original shifted variable by a union
bound. The other uses even convexity of the midpoint increment directly
on the difference.

Section~\ref{sec:conditional} establishes the finite-rank projections,
the two coefficient estimates, and the common passage to completed tails
and arbitrary centers. Section~\ref{sec:exponential} treats the exponential
law. Section~\ref{sec:gaussian} gives the two Gaussian routes.
Section~\ref{sec:obstruction} proves the renorming obstruction locally:
path products have norm one, but their child increments are weakly null
with fixed positive norm. An AUC renorming would force geometric growth
along a selected branch, contradicting boundedness.
